% TOP_PROBLEM: {"id": 3319, "title": "The Crossing Number of the Complete Bipartite Graph", "category_id": 3, "category": {"id": 3, "name": "graph_theory", "display_name": "Graph Theory", "description": "Problems involving graphs, networks, and their properties.", "slug": "graph-theory", "order_index": 3, "created_at": "2026-07-31T15:26:25.671Z"}, "status": "open", "problem_number": "OPG-310", "set_id": 12}
% FIRST_POSTED: 2026-10-01
% ENTRY_KIND: statement_only
% UnsolvedMath ID 3319; source catalogue code OPG-310
% !TeX program = lualatex
% Definition and sources only; this file is not an LLM solution attempt.
\documentclass[11pt,a4paper]{article}
\usepackage[margin=25mm]{geometry}
\usepackage{fontspec}
\setmainfont{lmroman10-regular.otf}[BoldFont=lmroman10-bold.otf,
 ItalicFont=lmroman10-italic.otf,BoldItalicFont=lmroman10-bolditalic.otf]
\usepackage{amsmath,amssymb,mathtools}
\usepackage{unicode-math}
\setmathfont{latinmodern-math.otf}
\AtBeginDocument{\let\setminus\smallsetminus}
\usepackage{xurl,hyperref}
\hypersetup{colorlinks=true,urlcolor=blue}
\setcounter{secnumdepth}{0}
\setlength{\parindent}{0pt}
\setlength{\parskip}{5pt}
\setlength{\emergencystretch}{3em}
\begin{document}
\section*{The Crossing Number of the Complete Bipartite Graph}\label{3319}
{\small UnsolvedMath ID 3319 · OPG-310 · Graph Theory · OpenGarden}
\subsection{Definitions and mathematical statement}
For integers $m,n\ge1$, the complete bipartite graph $K_{m,n}$ has
two disjoint independent vertex sets of sizes $m$ and $n$, with all
$mn$ possible edges between them. In a plane drawing, vertices are
distinct points and each edge is a simple arc between its ends,
avoiding other vertices. Intersections of edge interiors are finite
transverse crossings, with no three edges meeting at one crossing.
Endpoints shared by adjacent edges are not counted.

The crossing number $\operatorname{cr}(G)$ minimizes the number of
crossing points over all such topological drawings. The conjecture
is that for every pair of positive integers $m,n$,
\[
 \operatorname{cr}(K_{m,n})=
       \left\lfloor\frac m2\right\rfloor
       \left\lfloor\frac{m-1}2\right\rfloor
       \left\lfloor\frac n2\right\rfloor
       \left\lfloor\frac{n-1}2\right\rfloor.
\]
Floors round down. The expression is symmetric in the two part sizes
and vanishes when either part has at most two vertices.

A standard construction places the two parts on the two coordinate
axes, splitting each part as evenly as possible between the positive
and negative rays, and joins opposite-part pairs. Suitable generic
positions avoid triple crossings and realize the displayed count.
The target is the matching lower bound for every drawing.

There is no requirement that the two vertex parts occupy separate
regions, that all vertices lie on one curve, or that edges be
straight. Allowing arbitrary edge arcs is part of the minimization.
The conjecture is an exact finite formula, so an asymptotic estimate
with the same leading term would not alone determine it.
\subsection{Short English statement}
Does the complete bipartite graph always have crossing number equal to the product of the four displayed floor factors?
\subsection{Sources}
\begin{itemize}
\item[S1] ULAM.AI and the UnsolvedMath Contributors, supplied problems.json, record 3319 (OPG-310), OpenGarden collection. Catalogue identity and source transcription; CC BY 4.0.
\url{https://huggingface.co/datasets/ulamai/UnsolvedMath}.
\item[S2] Open Problem Garden, The Crossing Number of the Complete Bipartite Graph. Original problem and discussion; the mathematical formulation below is expanded from the supplied source record.
\url{http://www.openproblemgarden.org/op/the_crossing_number_of_the_complete_bipartite_graph}.
\end{itemize}
\end{document}
